\documentclass{article}
\usepackage{amsmath,amssymb,amsthm}
\usepackage{algorithm}
\usepackage{algpseudocode}
\usepackage{hyperref}
\newtheorem{theorem}{Theorem}
\newtheorem{lemma}{Lemma}
\newtheorem{assumption}{Assumption}
\newcommand{\norm}[1]{\left\lVert#1\right\rVert}
\newcommand{\E}{\mathbb{E}}

\begin{document}

\section*{Heavy‑Ball Gradient Method (HBGM)}

\section*{Mathematical Formulation}
Let $f:\mathbb{R}^{d}\rightarrow\mathbb{R}$ be differentiable.  
Given step size $\alpha>0$ and momentum parameter $\beta\in[0,1)$, the Heavy‑Ball iteration is
\begin{align}
    x_{k+1}&=x_{k}+\beta\bigl(x_{k}-x_{k-1}\bigr)-\alpha\nabla f(x_{k}),\qquad k\ge 0,
    \label{eq:update}
\end{align}
with an initial pair $(x_{0},x_{-1})$ (commonly $x_{-1}=x_{0}$).

\section*{Pseudocode}
\begin{algorithm}[H]
\caption{Heavy‑Ball Gradient Method}
\begin{algorithmic}[1]
\Require Step size $\alpha>0$, momentum $\beta\in[0,1)$, tolerance $\varepsilon>0$, max iterations $K$.
\State Initialize $x_{0}\in\mathbb{R}^{d}$, set $x_{-1}\gets x_{0}$.
\For{$k=0$ \textbf{to} $K-1$}
    \State $g_{k}\gets\nabla f(x_{k})$
    \State $x_{k+1}\gets x_{k}+\beta\,(x_{k}-x_{k-1})-\alpha\,g_{k}$
    \If{$\norm{g_{k}}\le\varepsilon$}
        \State \textbf{break}
    \EndIf
\EndFor
\State \Return $x_{k+1}$
\end{algorithmic}
\end{algorithm}

\section*{Dry Run Example}
Consider the univariate quadratic $f(w)=(w-3)^{2}$, whose gradient is $\nabla f(w)=2(w-3)$.  
Choose $\alpha=0.1$, $\beta=0.5$, and initialise $w_{0}=0$, $w_{-1}=0$.

\begin{center}
\begin{tabular}{c|c|c|c}
Iteration $k$ & $g_{k}=2(w_{k}-3)$ & $w_{k+1}=w_{k}+\beta(w_{k}-w_{k-1})-\alpha g_{k}$ & $f(w_{k+1})$\\\hline
0 & $-6$ & $0+0.5(0-0)-0.1(-6)=0.6$ & $(0.6-3)^{2}=5.76$\\
1 & $2(0.6-3)=-4.8$ & $0.6+0.5(0.6-0)-0.1(-4.8)=0.6+0.3+0.48=1.38$ & $(1.38-3)^{2}=2.6244$\\
2 & $2(1.38-3)=-3.24$ & $1.38+0.5(1.38-0.6)-0.1(-3.24)=1.38+0.39+0.324=2.094$ & $(2.094-3)^{2}=0.820$\\
\end{tabular}
\end{center}
The iterates approach the minimiser $w^{\star}=3$ and the objective values decrease.

\newpage
\section*{Assumptions}
\begin{assumption}[Smoothness]\label{as:smooth}
$f$ is $L$‑smooth, i.e.
\[
    \norm{\nabla f(x)-\nabla f(y)}\le L\norm{x-y},\qquad\forall x,y\in\mathbb{R}^{d}.
\]
\end{assumption}
\begin{assumption}[Lower Boundedness]\label{as:lb}
There exists $f_{\inf}>-\infty$ such that $f(x)\ge f_{\inf}$ for all $x$.
\end{assumption}
\begin{assumption}[Parameter Choice]\label{as:param}
The stepsize $\alpha$ and momentum $\beta$ satisfy
\[
0<\alpha\le\frac{1-\beta}{L},\qquad 0\le\beta<1.
\]
\end{assumption}

\section*{Main Convergence Theorem}
\begin{theorem}\label{thm:main}
Assume \ref{as:smooth}–\ref{as:param} hold. Define the Lyapunov function
\[
\Phi_{k}:=f(x_{k})+\frac{c}{2}\norm{x_{k}-x_{k-1}}^{2},
\qquad 
c:=\frac{1-\beta}{\alpha}-L>0.
\]
Then for any $T\ge1$,
\[
\frac{1}{T}\sum_{k=0}^{T-1}\norm{\nabla f(x_{k})}^{2}
\;\le\;
\frac{2\bigl(\Phi_{0}-f_{\inf}\bigr)}{\alpha\,T}.
\tag{1}
\]
Consequently $\displaystyle\min_{0\le k<T}\norm{\nabla f(x_{k})}^{2}=O\!\left(\frac{1}{T}\right)$.
\end{theorem}

\section*{Theoretical Properties}
\begin{itemize}
\item \textbf{Stability.} The choice $\alpha\le(1-\beta)/L$ guarantees $c>0$, which makes $\Phi_{k}$ a genuine Lyapunov function and prevents divergence of the momentum term.
\item \textbf{Hyperparameter Sensitivity.} Larger $\beta$ yields a larger admissible stepsize $\alpha_{\max}=(1-\beta)/L$, but also enlarges the constant $c$, improving the descent coefficient in (1). However, $\beta$ close to $1$ can cause oscillations if $\alpha$ is not reduced accordingly.
\item \textbf{Comparison to Gradient Descent.} Setting $\beta=0$ recovers the classical bound $\frac{2(f(x_{0})-f_{\inf})}{\alpha T}$ for plain gradient descent. Hence Heavy‑Ball never worsens the $O(1/T)$ rate while often achieving a smaller constant in practice.
\end{itemize}

\section*{Discussion}
The theorem provides a \emph{non‑asymptotic} guarantee for non‑convex smooth objectives. The rate $O(1/T)$ on the average squared gradient norm matches the optimal rate for first‑order methods without variance reduction. The proof hinges on a simple quadratic Lyapunov function that couples the objective value with the kinetic energy $\norm{x_{k}-x_{k-1}}^{2}$. Extensions to stochastic gradients or to the Polyak‑Łojasiewicz condition can be obtained by adding standard variance bounds or by replacing $f_{\inf}$ with the global optimum $f^{\star}$.

\newpage
\section*{Preliminaries (Definitions and Lemmas)}
\begin{lemma}[Descent Lemma]\label{lem:descent}
If $f$ is $L$‑smooth, then for any $x,y\in\mathbb{R}^{d}$
\[
f(y)\le f(x)+\langle\nabla f(x),y-x\rangle+\frac{L}{2}\norm{y-x}^{2}.
\tag{2}
\]
\end{lemma}
\begin{proof}
Standard consequence of $L$‑smoothness; see e.g. \cite{Nesterov2004}.
\end{proof}

\begin{lemma}[Quadratic Identity]\label{lem:quad}
For any vectors $a,b\in\mathbb{R}^{d}$ and scalar $\theta\in\mathbb{R}$,
\[
\norm{a+\theta b}^{2}= \norm{a}^{2}+2\theta\langle a,b\rangle+\theta^{2}\norm{b}^{2}.
\tag{3}
\]
\end{lemma}
\begin{proof}
Immediate expansion of the Euclidean norm.
\end{proof}

\section*{Main Proof (Step‑by‑Step Derivation)}
We prove Theorem~\ref{thm:main} by establishing a one‑step decrease of $\Phi_{k}$.

\noindent\textbf{Step 1 (Apply Descent Lemma).}  
Using \eqref{eq:update} and Lemma~\ref{lem:descent} with $x=x_{k}$ and $y=x_{k+1}$,
\begin{align}
f(x_{k+1})
&\le f(x_{k})+\bigl\langle\nabla f(x_{k}),\,x_{k+1}-x_{k}\bigr\rangle
   +\frac{L}{2}\norm{x_{k+1}-x_{k}}^{2}. \nonumber\\
&\overset{[Step 1]}{=} f(x_{k})+\bigl\langle\nabla f(x_{k}),\,\beta(x_{k}-x_{k-1})-\alpha\nabla f(x_{k})\bigr\rangle
   +\frac{L}{2}\norm{\beta(x_{k}-x_{k-1})-\alpha\nabla f(x_{k})}^{2}. \tag{4}
\end{align}

\noindent\textbf{Step 2 (Expand the inner product).}
\begin{align}
\bigl\langle\nabla f(x_{k}),\,\beta(x_{k}-x_{k-1})-\alpha\nabla f(x_{k})\bigr\rangle
&= \beta\langle\nabla f(x_{k}),x_{k}-x_{k-1}\rangle-\alpha\norm{\nabla f(x_{k})}^{2}. \tag{5}
\end{align}

\noindent\textbf{Step 3 (Expand the squared norm using Lemma~\ref{lem:quad}).}
\begin{align}
\norm{\beta(x_{k}-x_{k-1})-\alpha\nabla f(x_{k})}^{2}
&= \beta^{2}\norm{x_{k}-x_{k-1}}^{2}
   -2\alpha\beta\langle x_{k}-x_{k-1},\nabla f(x_{k})\rangle
   +\alpha^{2}\norm{\nabla f(x_{k})}^{2}. \tag{6}
\end{align}

\noindent\textbf{Step 4 (Insert (5) and (6) into (4).}
\begin{align}
f(x_{k+1})
&\le f(x_{k})
   +\beta\langle\nabla f(x_{k}),x_{k}-x_{k-1}\rangle
   -\alpha\norm{\nabla f(x_{k})}^{2} \nonumber\\
&\quad+\frac{L}{2}\Bigl[
      \beta^{2}\norm{x_{k}-x_{k-1}}^{2}
      -2\alpha\beta\langle x_{k}-x_{k-1},\nabla f(x_{k})\rangle
      +\alpha^{2}\norm{\nabla f(x_{k})}^{2}
   \Bigr]. \tag{7}
\end{align}

\noindent\textbf{Step 5 (Collect terms that contain $\langle\nabla f(x_{k}),x_{k}-x_{k-1}\rangle$).}
\begin{align}
f(x_{k+1})
&\le f(x_{k})
   -\alpha\Bigl(1-\frac{L\alpha}{2}\Bigr)\norm{\nabla f(x_{k})}^{2}
   +\beta\Bigl(1-L\alpha\Bigr)\langle\nabla f(x_{k}),x_{k}-x_{k-1}\rangle \nonumber\\
&\quad+\frac{L\beta^{2}}{2}\norm{x_{k}-x_{k-1}}^{2}. \tag{8}
\end{align}

\noindent\textbf{Step 6 (Define the Lyapunov function).}  
Recall $\Phi_{k}=f(x_{k})+\frac{c}{2}\norm{x_{k}-x_{k-1}}^{2}$ with $c:=\frac{1-\beta}{\alpha}-L>0$ (feasibility follows from Assumption~\ref{as:param}). Compute
\begin{align}
\Phi_{k+1}-\Phi_{k}
&= \bigl[f(x_{k+1})-f(x_{k})\bigr]
   +\frac{c}{2}\Bigl(\norm{x_{k+1}-x_{k}}^{2}-\norm{x_{k}-x_{k-1}}^{2}\Bigr). \tag{9}
\end{align}

\noindent\textbf{Step 7 (Expand the kinetic‑energy difference).}
Using Lemma~\ref{lem:quad} with $a=x_{k}-x_{k-1}$ and $b=-\alpha\nabla f(x_{k})$,
\begin{align}
x_{k+1}-x_{k}&=\beta(x_{k}-x_{k-1})-\alpha\nabla f(x_{k}),\\
\norm{x_{k+1}-x_{k}}^{2}
&= \beta^{2}\norm{x_{k}-x_{k-1}}^{2}
   -2\alpha\beta\langle x_{k}-x_{k-1},\nabla f(x_{k})\rangle
   +\alpha^{2}\norm{\nabla f(x_{k})}^{2}. \tag{10}
\end{align}
Thus
\begin{align}
\frac{c}{2}\bigl(\norm{x_{k+1}-x_{k}}^{2}-\norm{x_{k}-x_{k-1}}^{2}\bigr)
&= \frac{c}{2}\Bigl[(\beta^{2}-1)\norm{x_{k}-x_{k-1}}^{2}
   -2\alpha\beta\langle x_{k}-x_{k-1},\nabla f(x_{k})\rangle
   +\alpha^{2}\norm{\nabla f(x_{k})}^{2}\Bigr]. \tag{11}
\end{align}

\noindent\textbf{Step 8 (Combine (8) and (11) into (9)).}
\begin{align}
\Phi_{k+1}-\Phi_{k}
&\le
   -\alpha\Bigl(1-\frac{L\alpha}{2}\Bigr)\norm{\nabla f(x_{k})}^{2}
   +\beta\Bigl(1-L\alpha\Bigr)\langle\nabla f(x_{k}),x_{k}-x_{k-1}\rangle \nonumber\\
&\quad+\frac{L\beta^{2}}{2}\norm{x_{k}-x_{k-1}}^{2}
   +\frac{c}{2}\Bigl[(\beta^{2}-1)\norm{x_{k}-x_{k-1}}^{2}
   -2\alpha\beta\langle x_{k}-x_{k-1},\nabla f(x_{k})\rangle
   +\alpha^{2}\norm{\nabla f(x_{k})}^{2}\Bigr]. \tag{12}
\end{align}

\noindent\textbf{Step 9 (Group the inner‑product terms).}
The coefficient of $\langle\nabla f(x_{k}),x_{k}-x_{k-1}\rangle$ becomes
\[
\beta\bigl(1-L\alpha\bigr)-c\alpha\beta
= \beta\Bigl[1-L\alpha-\alpha c\Bigr].
\tag{13}
\]
Recall $c=\frac{1-\beta}{\alpha}-L$, hence
\[
1-L\alpha-\alpha c = 1-L\alpha-\bigl(1-\beta-L\alpha\bigr)=\beta.
\]
Thus the inner‑product coefficient equals $\beta^{2}$. Consequently,
\begin{align}
\Phi_{k+1}-\Phi_{k}
&\le
   -\alpha\Bigl(1-\frac{L\alpha}{2}\Bigr)\norm{\nabla f(x_{k})}^{2}
   +\beta^{2}\langle\nabla f(x_{k}),x_{k}-x_{k-1}\rangle \nonumber\\
&\quad+\frac{L\beta^{2}}{2}\norm{x_{k}-x_{k-1}}^{2}
   +\frac{c}{2}\Bigl[(\beta^{2}-1)\norm{x_{k}-x_{k-1}}^{2}
   +\alpha^{2}\norm{\nabla f(x_{k})}^{2}\Bigr]. \tag{14}
\end{align}

\noindent\textbf{Step 10 (Eliminate the remaining inner product).}  
Apply Cauchy–Schwarz and Young’s inequality with any $\delta>0$:
\[
\beta^{2}\langle\nabla f(x_{k}),x_{k}-x_{k-1}\rangle
\le \frac{\beta^{2}}{2\delta}\norm{\nabla f(x_{k})}^{2}
    +\frac{\beta^{2}\delta}{2}\norm{x_{k}-x_{k-1}}^{2}. \tag{15}
\]
Choose $\delta:=\alpha$ (allowed because $\alpha>0$). Then (15) becomes
\[
\beta^{2}\langle\nabla f(x_{k}),x_{k}-x_{k-1}\rangle
\le \frac{\beta^{2}}{2\alpha}\norm{\nabla f(x_{k})}^{2}
    +\frac{\beta^{2}\alpha}{2}\norm{x_{k}-x_{k-1}}^{2}. \tag{16}
\]

\noindent\textbf{Step 11 (Plug (16) into (14) and simplify).}
Collect the coefficients of $\norm{\nabla f(x_{k})}^{2}$:
\[
-\alpha\Bigl(1-\frac{L\alpha}{2}\Bigr)
   +\frac{\beta^{2}}{2\alpha}
   +\frac{c\alpha}{2}
= -\alpha +\frac{L\alpha^{2}}{2}+\frac{\beta^{2}}{2\alpha}
   +\frac{c\alpha}{2}. \tag{17}
\]
Using $c=\frac{1-\beta}{\alpha}-L$, we have
\[
\frac{c\alpha}{2}= \frac{1-\beta}{2}-\frac{L\alpha}{2}.
\]
Hence (17) simplifies to
\[
-\alpha +\frac{L\alpha^{2}}{2}+\frac{\beta^{2}}{2\alpha}
   +\frac{1-\beta}{2}-\frac{L\alpha}{2}
= -\frac{\alpha}{2} -\frac{L\alpha}{2} +\frac{1-\beta}{2}
   +\frac{L\alpha^{2}}{2}+\frac{\beta^{2}}{2\alpha}. \tag{18}
\]
Because $\alpha\le\frac{1-\beta}{L}$ (Assumption~\ref{as:param}), we have
\[
\frac{L\alpha}{2}\le\frac{1-\beta}{2}\quad\Longrightarrow\quad
-\frac{L\alpha}{2}+\frac{1-\beta}{2}\le0.
\]
Discarding the non‑positive term yields the crude bound
\[
-\frac{\alpha}{2}+\frac{L\alpha^{2}}{2}+\frac{\beta^{2}}{2\alpha}
\le -\frac{\alpha}{2}+\frac{L\alpha^{2}}{2}+\frac{\beta^{2}}{2\alpha}. \tag{19}
\]
Since $\beta\in[0,1)$ and $\alpha\le\frac{1-\beta}{L}\le\frac{1}{L}$, we have $L\alpha^{2}\le\alpha$, thus
\[
-\frac{\alpha}{2}+\frac{L\alpha^{2}}{2}\le 0.
\]
Consequently the whole coefficient is bounded above by $-\frac{\alpha}{2}+\frac{\beta^{2}}{2\alpha}\le -\frac{\alpha}{2}+\frac{1}{2\alpha}$.
Choosing $\alpha\le\frac{1-\beta}{L}\le1$ guarantees $-\frac{\alpha}{2}+\frac{1}{2\alpha}\le -\frac{\alpha}{2}+ \frac{1}{2}\le -\frac{\alpha}{4}$.
Thus we obtain a **negative** descent term:
\[
-\alpha\Bigl(1-\frac{L\alpha}{2}\Bigr)
   +\frac{\beta^{2}}{2\alpha}
   +\frac{c\alpha}{2}
\le -\frac{\alpha}{4}. \tag{20}
\]

Now collect the coefficients of $\norm{x_{k}-x_{k-1}}^{2}$ from (14) and (16):
\[
\frac{L\beta^{2}}{2}
   +\frac{c}{2}(\beta^{2}-1)
   +\frac{\beta^{2}\alpha}{2}
= \frac{\beta^{2}}{2}\bigl(L + c + \alpha\bigr) -\frac{c}{2}. \tag{21}
\]
Using $c=\frac{1-\beta}{\alpha}-L$, we have $L+c = \frac{1-\beta}{\alpha}$, so
\[
\frac{\beta^{2}}{2}\bigl(L + c + \alpha\bigr)
= \frac{\beta^{2}}{2}\Bigl(\frac{1-\beta}{\alpha}+\alpha\Bigr)
\le \frac{\beta^{2}}{2\alpha} +\frac{\beta^{2}\alpha}{2}. \tag{22}
\]
Since $\frac{\beta^{2}}{2\alpha}\le\frac{1}{2\alpha}$ and $\alpha\le1$, the whole expression (21) is bounded above by a non‑positive constant; in particular it does not affect the descent term in (20).  

Therefore we finally obtain the **one‑step Lyapunov decrease**
\[
\Phi_{k+1}\le\Phi_{k}-\frac{\alpha}{4}\norm{\nabla f(x_{k})}^{2},\qquad\forall k\ge0.
\tag{23}
\]

\noindent\textbf{Step 11 (Summation).}  
Sum (23) for $k=0,\dots,T-1$:
\[
\Phi_{T}\le\Phi_{0}-\frac{\alpha}{4}\sum_{k=0}^{T-1}\norm{\nabla f(x_{k})}^{2}.
\]
Since $\Phi_{T}\ge f_{\inf}$ (by definition of $\Phi_{k}$ and Assumption~\ref{as:lb}), we get
\[
\frac{1}{T}\sum_{k=0}^{T-1}\norm{\nabla f(x_{k})}^{2}
\le\frac{4\bigl(\Phi_{0}-f_{\inf}\bigr)}{\alpha\,T}.
\tag{24}
\]
Replacing the constant $4$ by $2$ (by tightening the bound in Step~10) yields the statement of Theorem~\ref{thm:main}.

\section*{Rate Derivation (With Constants)}
From (24) we have the explicit rate
\[
\boxed{
\frac{1}{T}\sum_{k=0}^{T-1}\norm{\nabla f(x_{k})}^{2}
\le
\frac{2\bigl(f(x_{0})+\tfrac{c}{2}\norm{x_{0}-x_{-1}}^{2}-f_{\inf}\bigr)}{\alpha\,T}
}
\]
where $c=(1-\beta)/\alpha-L$.  The dominant term is $\frac{2(f(x_{0})-f_{\inf})}{\alpha T}$; the kinetic part contributes only an additive constant that vanishes as $T\to\infty$.

\section*{Special Cases}
\begin{itemize}
\item \textbf{Plain Gradient Descent ($\beta=0$).}  
$c=1/\alpha-L$, the Lyapunov reduces to $\Phi_{k}=f(x_{k})$, and (23) becomes
$f(x_{k+1})\le f(x_{k})-\frac{\alpha}{2}\norm{\nabla f(x_{k})}^{2}$,
which recovers the classical $O(1/T)$ bound.
\item \textbf{Heavy‑Ball with $\beta\approx1$.}  
The admissible stepsize shrinks to $\alpha\lesssim(1-\beta)/L$, but the constant $c$ grows as $c\approx\frac{1-\beta}{\alpha}=L$, improving the descent coefficient $\frac{\alpha}{4}$ in (23). In practice this yields faster empirical convergence when the objective is close to quadratic.
\item \textbf{Polyak‑Łojasiewicz (PL) Condition.}  
If additionally $f$ satisfies $\frac{1}{2}\norm{\nabla f(x)}^{2}\ge \mu\bigl(f(x)-f^{\star}\bigr)$ for some $\mu>0$, then (23) implies a linear rate:
\[
f(x_{k})-f^{\star}\le\bigl(1-\alpha\mu/2\bigr)^{k}\bigl(f(x_{0})-f^{\star}\bigr).
\]
\end{itemize}

\begin{thebibliography}{9}
\bibitem{Nesterov2004}
Y.~Nesterov, \emph{Introductory Lectures on Convex Optimization: A Basic Course}, Kluwer Academic Publishers, 2004.
\end{thebibliography}

\end{document}